1.7. 若 $Y$ 为域 $k$ 上的射影空间 $\mathbb P^r(k)$ 中维数至少为三的完全交, 则已知 $\operatorname{Pic}(Y)$ 同构于 $\mathbb Z$, 并由 $\underline O(1)$ 生成 (SGA 2, XII, 3.7). 在二维时, 若 $Y$ 光滑, 仍可证明如下结果.

\underline{定理} 1.8. \underline{设 $Y$ 为代数闭域 $k$ 上射影空间 $\mathbb P^r(k)$ 中的光滑完全交曲面. 则 $\underline{\operatorname{Pic}}^\tau(Y)=0$, 且 $\underline O(1)$ 的类在 Néron--Severi 群 $\operatorname{NS}(Y)$ 中不可除.}

由 1.5(iii)$_a$ 知 $H^1(Y,\underline O)=0$, 因而 $\underline{\operatorname{Pic}}^0(Y)=0$. 再由下面的熟知引理、1.6(iv) 所保证的 $H^2(Y,\mathbb Z_\ell(1))$ 无挠性, 以及 1.5(iii)$_a$ 给出的 $H^0(Y,\Omega_Y^1)=0$, 即得 $\underline{\operatorname{Pic}}^\tau(Y)$ 为零.

\underline{引理} 1.9. \underline{设 $Y$ 为代数闭域 $k$ 上的完备代数簇.}

\begin{enumerate}[label=(\roman*),leftmargin=2.6em]
\item \underline{若 $\ell$ 与 $k$ 的特征互素, 则 $\operatorname{NS}(Y)$ 的 $\ell$ 挠子群可等同于 $H^2(Y,\mathbb Z_\ell(1))$ 的 $\ell$ 挠子群.}

%%%%%%%%%% source scan idx 57 | folio 50 %%%%%%%%%%

\item \underline{若 $k$ 的特征 $p$ 非零且 $Y$ 光滑, 则 $\operatorname{Pic}(Y)$ 中乘以 $p$ 的核可等同于 $H^0(Y,\Omega_Y^1)$ 中由局部对数微分组成的子群.}
\end{enumerate}

Kummer 正合列

\begin{equation*}
0\longrightarrow\mu_{\ell^n}\longrightarrow\mathbb G_m
\xrightarrow{\ x\mapsto x^{\ell^n}\ }\mathbb G_m\longrightarrow0
\end{equation*}

给出正合列

\begin{equation*}
0\longrightarrow
\operatorname{Pic}(Y)/\ell^n\operatorname{Pic}(Y)
\longrightarrow H^2(Y,\mu_{\ell^n})
\longrightarrow H^2(Y,\mathbb G_m)_{\ell^n}
\longrightarrow0,
\end{equation*}

其中

\begin{equation*}
\operatorname{Pic}(Y)/\ell^n\operatorname{Pic}(Y)
=\operatorname{NS}(Y)\otimes\mathbb Z/\ell^n.
\end{equation*}

取逆极限得到

\begin{equation*}
0\longrightarrow\operatorname{NS}(Y)\otimes\mathbb Z_\ell
\longrightarrow H^2(Y,\mathbb Z_\ell(1))
\longrightarrow
\operatorname{Hom}\bigl(\mathbb Q_\ell/\mathbb Z_\ell,
H^2(Y,\mathbb G_m)\bigr)
\longrightarrow0.
\end{equation*}

由于 $\mathbb Q_\ell/\mathbb Z_\ell$ 可除, 余核无挠. 这就证明了 (i).

为证明 (ii), 只须注意正合列

\begin{equation*}
0\longrightarrow\underline O_Y^*
\xrightarrow{\ x\mapsto x^p\ }\underline O_Y^*
\xrightarrow{\ df/f\ }\Omega_Y^1
\end{equation*}

定义了短正合列

\begin{equation*}
0\longrightarrow\underline O_Y^*\longrightarrow\underline O_Y^*
\longrightarrow\Omega^1_{Y,\mathrm{loc}\,\log}\longrightarrow0,
\end{equation*}

再取相应的上同调长正合列即可. 此论证出自 Cartier; 他还证明, 若以 $C$ 表示 $H^0(Y,\Omega_Y^1)$ 上的 Cartier 算子, 则局部对数微分构成 $C-\operatorname{Id}$ 的核.

反设 $\underline O(1)$ 在 $\operatorname{NS}(Y)$ 中的类具有形式 $nu$, 其中 $n>1$. 若 $\ell$ 整除 $n$ 且与特征互素, 则在 $\ell$ 进上同调中仍有

%%%%%%%%%% source scan idx 58 | folio 51 %%%%%%%%%%

\begin{equation*}
\eta_\ell=n\,c_1(u),
\end{equation*}

这与 1.6(iv) 矛盾.

最后, 若 $k$ 的特征 $p$ 非零且 $p$ 整除 $n$, 则在 Hodge 上同调中仍有

\begin{equation*}
\eta=n\,c_1(u)=0,
\end{equation*}

这与 1.5(iii)$_d$ 矛盾, 证明完毕.

